\clearpage\fancyhead[L]{\small COMET.}\hypertarget{native-394-lower}{}\label{native-394-lower}
\begin{center}\textbf{COMET.}\par\vspace{4pt}\hrule\vspace{4pt}\fontsize{10}{12}\selectfont Göttingische gelehrte Anzeigen. Stück 149. Seite 1473..1476. 1815 Septemb. 18.\par\vspace{4pt}\hrule\vspace{4pt}\end{center}

\par\indent Die Beobachtungen des so höchst merkwürdigen Cometen, von welchem wir mehrere Male in
diesen Blättern Nachrichten mitgetheilt haben, sind auf der hiesigen Sternwarte von dem Herrn Prof.
\textsc{Gauss} länger fortgesetzt worden, als sich anfangs hoffen liess. Da sie gegenwärtig geschlossen sind, so
stellen wir sie hier vollständig zusammen.

\SourceLeaf{395}\begin{center}\fontsize{10}{12}\selectfont\setlength{\tabcolsep}{5pt}\renewcommand{\arraystretch}{1.15}\begin{tabular}{l|c|c|c}
1815 & \shortstack{Mittlere Zeit\\in Göttingen} & \shortstack{Scheinb. gerade\\Aufsteigung} & \shortstack{Scheinbare\\Abweichung}\\\hline
März 20 & $10^{\mathrm u}\;33^{\mathrm m}\;6^{\mathrm s}$ & $54^\circ\;7'\;1''$ & $39^\circ\;7'\;47''\;\mathrm{N.}$\\
21 & $10\;11\;37$ & $54\;33\;21$ & $39\;36\;57$\\
25 & $9\;46\;38$ & $56\;28\;50$ & $41\;38\;5$\\
30 & $9\;50\;57$ & $59\;13\;3$ & $44\;10\;27$\\
April 2 & $9\;11\;55$ & $61\;2\;27$ & $45\;39\;57$\\
Juni 12 & $10\;53\;31$ & $165\;3\;3$ & $52\;8\;25$\\
30 & $10\;49\;12$ & $184\;39\;24$ & $39\;22\;5$\\
Juli 13 & $10\;57\;20$ & $194\;32\;52$ & $29\;53\;54$\\
27 & $10\;4\;59$ & $203\;0\;56$ & $20\;37\quad:$\\
29 & $10\;13\;58$ & $204\;6\;3$ & $19\;24\;13$\\
Aug. 4 & $10\;14\;7$ & $207\;14\;8$ & $15\;53\;49$\\
25 & $9\;1\;44$ & $217\;1\;31$ & $5\;33\;36$\\
\end{tabular}\end{center}

\par\indent Wir bemerken dabei noch, dass am 13. und 27. Juli und 4. August Sterne aus \textsc{Piazzis} Cataloge
zur Vergleichung angewandt sind, an allen übrigen Tagen hingegen Sterne, deren Positionen aus der
\emph{Histoire Céleste} oder den \emph{Mémoires de l'Ac. de Paris} 1790 reducirt sind. Die Beobachtungen sind alle
am Kreismikrometer des 10füssigen \textsc{Herschel}schen Teleskops angestellt, die vom 12. Juni ausgenommen,
wo der Abstand von 219 im grossen Bär nach \textsc{Bode} und der Richtungswinkel mit einem vortrefflichen
\textsc{Frauenhofer}schen Heliometer gemessen wurden.

\par\indent Die elliptische Bahn, welche Hr. Prof. \textsc{Gauss} um die Mitte des Juni bestimmt hatte, und die im
105. Stück dieser Blätter [S. 385 d. B.] mitgetheilt ist, hat auch mit allen spätern Beobachtungen eine
gute Uebereinstimmung beibehalten. Spätere und auf sorgfältigere Discussion der Beobachtungen gegründete Bestimmungen dieser Bahn, welche die Hrn. \textsc{Bessel} und \textsc{Nicolai} ausgeführt haben, weichen
gleichfalls nur wenig von jener ab, und die kurze Umlaufszeit dieses Cometen kann jetzt nur noch um
eine kleine Grösse ungewiss sein. Wir stellen hier diese Bestimmungen zusammen.

\Needspace{21\baselineskip}\begin{center}
\emph{Elliptische Elemente des Cometen berechnet vom Hrn. Prof.} \textsc{Bessel.}
\end{center}

\begin{center}\fontsize{10}{12}\selectfont Durchgang durch die Sonnennähe April 26.\par $0^{\mathrm u}\;5^{\mathrm m}\;14^{\mathrm s}5$ Mittlere Zeit in Paris\end{center}
\begin{center}\fontsize{9}{11}\selectfont\setlength{\tabcolsep}{2pt}\renewcommand{\arraystretch}{1.15}\begin{tabular}{lcr}
Länge des Knoten & $.\;.\;.\;.\;.\;.\;.\;.\;.\;.\;.\;.\;.$ & $83^\circ\;28'\;46''18$\\
Länge der Sonnennähe & $.\;.\;.\;.\;.\;.\;.\;.\;.\;.\;.\;.$ & $149\;2\;29.13$\\
\multicolumn{3}{c}{\shortstack{beide für den 1. Januar 1815 und\\siderisch ruhend vorausgesetzt}}\\
Neigung der Bahn & $.\;.\;.\;.\;.\;.\;.\;.\;.\;.\;.\;.\;.$ & $44\;29\;53.71$\\
Logarithm. des kleinsten Abstandes & $.\;.\;.\;.\;.\;.\;.$ & $0.0837950$\\
Excentricität & $.\;.\;.\;.\;.\;.\;.\;.\;.\;.\;.\;.\;.\;.\;.$ & $0.93112771$\\
Halbe grosse Axe & $.\;.\;.\;.\;.\;.\;.\;.\;.\;.\;.\;.\;.\;.$ & $17.60964$\\
Siderische Umlaufszeit & $.\;.\;.\;.\;.\;.\;.\;.\;.\;.\;.\;.$ & $73.89682\;\text{Jahre}$\\
\end{tabular}\end{center}

\Needspace{21\baselineskip}\begin{center}
\emph{Elliptische Elemente des Cometen, berechnet von Hrn.} \textsc{Nicolai}, \emph{auf der Seeberger Sternwarte.}
\end{center}

\begin{center}\fontsize{10}{12}\selectfont Durchgang durch die Sonnennähe April 26.\par $0^{\mathrm u}\;55^{\mathrm m}\;32^{\mathrm s}$ M. Z. der Seeberger Sternwarte\end{center}
\begin{center}\fontsize{9}{11}\selectfont\setlength{\tabcolsep}{2pt}\renewcommand{\arraystretch}{1.15}\begin{tabular}{lcr}
Länge der Sonnennähe & $.\;.\;.\;.\;.\;.\;.\;.\;.\;.\;.\;.$ & $149^\circ\;3'\;25''3$\\
Länge des Knoten & $.\;.\;.\;.\;.\;.\;.\;.\;.\;.\;.\;.\;.\;.$ & $83\;28\;52.3$\\
\multicolumn{3}{c}{beide für den 26. April und siderisch ruhend}\\
Neigung der Bahn & $.\;.\;.\;.\;.\;.\;.\;.\;.\;.\;.\;.\;.$ & $44\;29\;46.0$\\
Logarithm. des kleinsten Abstandes & $.\;.\;.\;.\;.\;.\;.$ & $0.0837490$\\
Excentricität & $.\;.\;.\;.\;.\;.\;.\;.\;.\;.\;.\;.\;.\;.\;.\;.$ & $0.93029341$\\
Halbe grosse Axe & $.\;.\;.\;.\;.\;.\;.\;.\;.\;.\;.\;.\;.\;.$ & $17.39704$\\
Siderische Umlaufszeit & $.\;.\;.\;.\;.\;.\;.\;.\;.\;.\;.\;.$ & $72.564\;\text{Julian. Jahre.}$\\
\end{tabular}\end{center}

\Needspace{13\baselineskip}\SourceLeaf{396}\begingroup\setlength{\rightskip}{0pt plus 1em}\par\indent Die Vergleichung der letztern Elemente mit den spätesten hiesigen Beobachtungen gab folgende\par\endgroup

\begin{center}
\fontsize{10}{12}\selectfont\setlength{\tabcolsep}{9pt}\renewcommand{\arraystretch}{1.15}\begin{tabular}{l | c | c}
& \multicolumn{2}{c}{Unterschiede} \\
& ger. Aufst. & Abweichung \\
\hline
Juni 30 & $-50''4$ & $+15''9$ \\
Juli 13 & $-16.6$ & $-44.1$ \\
27 & $-23.8$ & \\
29 & $+3.6$ & $-47.8$ \\
August 4 & $-6.5$ & $-38.1$ \\
25 & $-2.4$ & $-47.5$ \\
\end{tabular}
\end{center}

\Needspace{13\baselineskip}\par\indent Wir holen bei dieser Gelegenheit auch noch die Anzeige der Beobachtungen der \emph{Juno} nach,
welche vom Hrn. Prof. \textsc{Gauss} um die Zeit der Opposition gemacht sind. Es scheint nicht, dass dieser
diesmal ungemein lichtschwache Planet auf einer andern Sternwarte beobachtet wäre.

\begin{center}\fontsize{10}{12}\selectfont\setlength{\tabcolsep}{5pt}\renewcommand{\arraystretch}{1.15}\begin{tabular}{l|c|c|c}
1815 & \shortstack{Mittlere Zeit\\in Göttingen} & \shortstack{Scheinb. gerade\\Aufsteigung} & \shortstack{Scheinbare\\Abweichung}\\\hline
März 1 & $10^{\mathrm u}\;22^{\mathrm m}\;0''5$ & $197^\circ\;7'\;56''7$ & $2^\circ\;19'\;3''0\;\mathrm{S.}$\\
29 & $10\;8\;57.0$ & $192\;31\;51.8$ & $1\;32\;30.1\;\mathrm{N.}$\\
April 8 & $9\;27\;42.9$ & $190\;34\;51.9$ & $2\;53\;24.2\;\mathrm{N.}$\\
\end{tabular}\end{center}

\Needspace{14\baselineskip}\par\indent Hr. \textsc{Nicolai} übernahm die Vergleichung dieser Beobachtungen mit den letzten durch Hrn. \textsc{Möbius}
berechneten Elementen (S. Mon. Corresp. 1813. Bd. XXVIII. S. 577), und fand, nachdem er die Epoche
um $4'$ $55''$ vergrössert hatte, folgende

\begin{center}
\fontsize{10}{12}\selectfont\setlength{\tabcolsep}{6pt}\renewcommand{\arraystretch}{1.15}\begin{tabular}{l | c | c | c | c}
& \multicolumn{4}{c}{Unterschiede} \\
& ger. Aufst. & Abweich. & Länge & Breite \\
\hline
März 1 & $-20''6$ & $-49.7$ & $0$ & $-53''7$ \\
29 & $-12.2$ & $-50.5$ & $+8.6$ & $-51.2$ \\
April 8 & $-16.3$ & $-45.1$ & $+2.9$ & $-47.8$ \\
\end{tabular}
\end{center}

\noindent und hieraus die Opposition mit der Sonne

\begin{center}1815 März 31.\quad $12^{\mathrm u}\;56^{\mathrm m}\;43^{\mathrm s}$ Mittlere Zeit in Göttingen\end{center}
\begin{center}\fontsize{9}{11}\selectfont\setlength{\tabcolsep}{2pt}\renewcommand{\arraystretch}{1.15}\begin{tabular}{lcr}
Wahre Länge & $.\;.\;.\;.\;.\;.\;.\;.\;.\;.\;.\;.\;.$ & $190^\circ\;25'\;6''6$\\
Nördliche Geocentrische Breite & $.\;.\;.\;.\;.\;.\;.$ & $6\;29\;10.0$\\
\end{tabular}\end{center}

\Needspace{21\baselineskip}\par\indent Ferner erhielt dieser geschickte Astronom durch Verbindung dieser Opposition mit den drei vorhergehenden folgende

\Needspace{18\baselineskip}\begin{center}
\emph{Neue Elemente der Juno.}
\end{center}

\begin{center}\fontsize{9}{11}\selectfont\setlength{\tabcolsep}{2pt}\renewcommand{\arraystretch}{1.15}\begin{tabular}{lcr}
Epoche der mittlern Länge 1816 & $.\;.$ & $230^\circ\;11'\;34''2$\\
\multicolumn{3}{c}{für den Meridian von Göttingen}\\
Tägliche tropische Bewegung & $.\;.\;.$ & $812''9304$\\
Länge des Perihels 1816 & $.\;.\;.\;.\;.$ & $53^\circ\;14'\;53''8$\\
Länge des aufsteigenden Knoten 1816 &  & $171\;9\;58.9$\\
Neigung der Bahn & $.\;.\;.\;.\;.\;.\;.$ & $13\;14\;0.1$\\
Excentricitätswinkel & $.\;.\;.\;.\;.\;.$ & $14\;43\;28.84$\\
Logarithm. der halben grossen Axe & $.$ & $0.4266844$\\
\end{tabular}\end{center}

\par\indent Eine nach diesen Elementen von Hrn \textsc{Nicolai} berechnete Ephemeride wird in dem astronomischen Jahrbuche für 1818 bekannt gemacht werden.
